% ============================================================
%  PLANTILLA GUÍAS INSTITUCIONALES — USACH PREMIUM
%  Fusiona el estilo formal de guia1 (portada + marca de agua)
%  con el estilo premium de guia3 (colores + tcolorbox + multicol).
%
%  INSTRUCCIONES RÁPIDAS:
%  1. Edita el bloque "INFORMACIÓN DE LA GUÍA" (líneas ~65-85).
%  2. Rellena \listacontenidos con los temas de tu guía.
%  3. Agrega ejercicios en los entornos indicados.
%  4. Compila con: pdflatex guia-institucional.tex (2 veces).
%
%  ESTRUCTURA DEL DOCUMENTO:
%    Portada → Introducción → Kit de Propiedades → Ejercicios → Solucionario
% ============================================================

\documentclass[letterpaper, 12pt]{article}

% ============================================================
% PAQUETES
% ============================================================
\usepackage[utf8]{inputenc}
\usepackage[spanish]{babel}
\usepackage{amsmath, amssymb}
\usepackage{geometry}
\usepackage{fancyhdr}
\usepackage{enumitem}
\usepackage{graphicx}
\usepackage{hyperref}
\usepackage{multicol}
\usepackage{parskip}
\usepackage{xcolor}
\usepackage{tcolorbox}
\usepackage{eso-pic}
\usepackage{transparent}
\usepackage{tikz}
\usetikzlibrary{patterns}
\tcbuselibrary{skins}

% ============================================================
% GEOMETRÍA
% ============================================================
\geometry{letterpaper, margin=1.5cm, top=3cm, bottom=2.5cm, headheight=2cm}

% ============================================================
% COLORES INSTITUCIONALES
% ============================================================
\definecolor{celesteSuave}{HTML}{F0F7FF}
\definecolor{azulFuerte}{HTML}{1A5276}
\definecolor{turquesaUsach}{HTML}{33CCCC}
\definecolor{enlace}{HTML}{1A5276}

% ============================================================
% RUTAS DE IMÁGENES  ← ajusta si tu carpeta de imágenes es distinta
% ============================================================
\newcommand{\logoup}{./images/logo_up.png}
\newcommand{\logousach}{./images/logo_usach.png}
\newcommand{\logofondo}{./images/logo_up.png}

% ============================================================
% INFORMACIÓN DE LA GUÍA  ← EDITA AQUÍ ANTES DE COMPILAR
% ============================================================
\newcommand{\institucion}{Usach Premium}
\newcommand{\curso}{Nombre del Curso}              % ej: Cálculo I
\newcommand{\guiasnumero}{Guía N - Título}         % ej: Guía 1 - Trigonometría
\newcommand{\semestre}{1er. Semestre de 2026}
\newcommand{\preparadopor}{Usach Premium.}         % o el nombre del ayudante
\newcommand{\webinstitucion}{www.usachpremium.cl}
\newcommand{\instagraminstitucion}{https://www.instagram.com/usach.premium}
\newcommand{\transparencia}{0.05}                  % opacidad marca de agua (0.0–1.0)
\newcommand{\fraseportada}{"Por y para estudiantes"}

% Lista de contenidos que aparece en la portada dentro del tcolorbox.
% Agrega o quita \item según los temas de tu guía.
\newcommand{\listacontenidos}{
    \item[\color{azulFuerte}$\Rightarrow$] Contenido 1
    \item[\color{azulFuerte}$\Rightarrow$] Contenido 2
    \item[\color{azulFuerte}$\Rightarrow$] Contenido 3
    \item[\color{azulFuerte}$\Rightarrow$] Contenido 4
}

% ============================================================
% HYPERREF
% ============================================================
\hypersetup{
    colorlinks=true,
    linkcolor=enlace,
    urlcolor=enlace,
    citecolor=enlace,
    pdfborder={0 0 0},
}

% ============================================================
% ENCABEZADO Y PIE DE PÁGINA (páginas interiores)
% ============================================================
\pagestyle{fancy}
\fancyhf{}
\fancyhead[L]{%
    \includegraphics[height=1.2cm]{\logoup}%
}
\fancyhead[R]{%
    \begin{minipage}[b]{0.42\textwidth}
        \raggedleft
        \fontsize{9pt}{11pt}\selectfont
        \textbf{\color{azulFuerte}\institucion} \\
        \curso\ — \guiasnumero \\
        \textit{\color{gray}@usach.premium}
    \end{minipage}\hspace{0.1cm}%
    \includegraphics[height=1.2cm]{\logousach}%
}
\fancyfoot[L]{\fontsize{9pt}{11pt}\selectfont \textit{\fraseportada}}
\fancyfoot[R]{\fontsize{9pt}{11pt}\selectfont Página \thepage}
\renewcommand{\headrulewidth}{0.4pt}
\renewcommand{\footrulewidth}{0.4pt}

% Estilo portada: sin encabezado ni pie
\fancypagestyle{plain}{
    \fancyhf{}
    \renewcommand{\headrulewidth}{0pt}
    \renewcommand{\footrulewidth}{0pt}
}

% ============================================================
% MARCA DE AGUA (logo transparente al centro de cada página)
% Llama a \MarcaAgua dentro de un entorno para activarla.
% ============================================================
\newcommand{\MarcaAgua}{
    \AddToShipoutPicture{
        \AtPageCenter{
            \makebox(0,0){%
                \transparent{\transparencia}%
                \includegraphics[width=0.7\textwidth]{\logofondo}%
            }
        }
    }
}

% ============================================================
% CAJAS REUTILIZABLES
% ============================================================

% Caja de resumen / kit de propiedades
% Uso: \begin{kitpropiedades}{Título de la caja} ... \end{kitpropiedades}
\newenvironment{kitpropiedades}[1]{%
    \begin{tcolorbox}[colback=celesteSuave, colframe=azulFuerte,
        title=\textbf{#1}, fonttitle=\bfseries]
}{%
    \end{tcolorbox}
}

% Caja de desafío con sombra y título de color blanco sobre azul
% Uso: \begin{desafiobox}{Título} ... \end{desafiobox}
\newtcolorbox{desafiobox}[1]{
    colback=white,
    colframe=azulFuerte,
    fonttitle=\bfseries,
    coltitle=white,
    title=#1,
    arc=8pt,
    boxrule=1.2pt,
    enhanced,l
    drop shadow={black!5}
}

% Caja de bloque de soluciones
% Uso: \begin{solbox}{Título} ... \end{solbox}
\newenvironment{solbox}[1]{%
    \begin{tcolorbox}[colback=celesteSuave, colframe=azulFuerte,
        title=\textbf{#1}, fonttitle=\bfseries]
}{%
    \end{tcolorbox}
}

% ============================================================
% ENTORNOS DE SECCIÓN
% Cada entorno activa la marca de agua en las páginas que lo contienen.
% ============================================================

\newenvironment{introduccion}{%
    \section*{Introducción}%
    \MarcaAgua
}{}

\newenvironment{ejercicios}{%
    \MarcaAgua
}{}

\newenvironment{soluciones}{%
    \newpage
    \begin{center}
        \begin{tcolorbox}[colback=azulFuerte, colframe=azulFuerte,
            colupper=white, width=0.8\textwidth, center, arc=10pt]
            \centering\Large\textbf{SOLUCIONARIO}
        \end{tcolorbox}
    \end{center}
    \MarcaAgua
    \small
}{}

% ============================================================
\begin{document}
\thispagestyle{plain}

% ============================================================
% PORTADA
% ============================================================

% Logos en la portada: UP arriba a la izquierda, USACH arriba a la derecha
\AddToShipoutPicture*{%
    \put(54, 700){\includegraphics[width=2cm]{\logoup}}
    \put(420, 706){\includegraphics[width=5cm]{\logousach}}
}

\vspace*{0.5cm}

\begin{center}
    \color{azulFuerte}
    {\huge \textbf{GUÍA DE EJERCICIOS}} \\[0.5cm]
    {\Huge \textbf{\curso}} \\[0.3cm]
    {\Large \guiasnumero} \\[1.5cm]

    % Caja de contenidos
    \begin{tcolorbox}[colback=celesteSuave, colframe=azulFuerte, arc=10pt,
        width=0.85\textwidth, center, boxrule=1.5pt]
        \centering
        \vspace{0.2cm}
        {\Large \textbf{Contenidos}}
        \vspace{0.3cm}
        \color{black}
        \begin{itemize}[leftmargin=3cm]
            \listacontenidos
        \end{itemize}
        \vspace{0.2cm}
    \end{tcolorbox}
\end{center}

\vspace{1.5cm}

\begin{center}
    {\Large \textbf{\semestre}} \\[0.8cm]
    {\large \textbf{\preparadopor}}
\end{center}

\vspace{2cm}

% Pie de portada con links
\begin{center}
    {\color{azulFuerte}\hrule height 0.5pt}
    \vspace{0.3cm}
    \begin{minipage}{0.45\textwidth}
        \centering
        \textbf{Instagram:} \\
        \small\href{\instagraminstitucion}{@usach.premium}
    \end{minipage}
    \begin{minipage}{0.45\textwidth}
        \centering
        \textbf{Web:} \\
        \small\href{http://\webinstitucion}{\webinstitucion}
    \end{minipage}
\end{center}

\pagenumbering{arabic}
\newpage

% ============================================================
% INTRODUCCIÓN  (borra este bloque si no aplica)
% ============================================================
\begin{introduccion}
    Texto introductorio. Explica el propósito de la guía, qué
    conceptos se trabajan y por qué son importantes para el curso.

    % Imagen opcional (descomenta y ajusta la ruta):
    % \vspace{0.8cm}
    % \begin{center}
    %     \includegraphics[width=0.9\textwidth]{./images/IMAGEN.png}
    % \end{center}
\end{introduccion}

\newpage

% ============================================================
% KIT DE PROPIEDADES  (borra este bloque si no aplica)
% ============================================================
\begin{kitpropiedades}{RESUMEN DE CLASES: LO QUE DEBES SABER}
\begin{multicols}{2}
    \textbf{1. Primera categoría}
    \begin{itemize}[leftmargin=10pt]
        \item Propiedad A
        \item Propiedad B
    \end{itemize}

    \columnbreak

    \textbf{2. Segunda categoría}
    \begin{itemize}[leftmargin=10pt]
        \item Propiedad C
        \item Propiedad D
    \end{itemize}
\end{multicols}
\vspace{-0.2cm}
\begin{center}
    \small \textbf{Tip Pro:} Escribe aquí un tip útil para los estudiantes.
\end{center}
\end{kitpropiedades}

% ============================================================
% EJERCICIOS
% ============================================================
\begin{ejercicios}

% --- Sección con ejercicios cortos (4 columnas) ---
\subsection*{I. Nombre de la Primera Sección}
\begin{multicols}{4}
\small
\begin{enumerate}[label=\textbf{\arabic*.}, noitemsep, leftmargin=*]
    \item Ejercicio corto 1
    \item Ejercicio corto 2
    \item Ejercicio corto 3
    \item Ejercicio corto 4
\end{enumerate}
\end{multicols}

\vspace{0.5cm}

% --- Sección con ejercicios medianos (2 columnas) ---
\subsection*{II. Nombre de la Segunda Sección}
\begin{multicols}{2}
\begin{enumerate}[label=\textbf{\arabic*.}, start=5, leftmargin=*]
    \item Ejercicio más largo 1
    \item Ejercicio más largo 2
\end{enumerate}
\end{multicols}

\vspace{0.5cm}

% --- Sección con sub-ítems (1 columna) ---
\subsection*{III. Nombre de la Tercera Sección}
\begin{enumerate}[label=\arabic*., leftmargin=*, itemsep=0.4cm]
    \item Enunciado principal con sub-partes:
    \begin{enumerate}[label=\arabic*), leftmargin=*, itemsep=0.4cm]
        \item $\displaystyle $ % Escribe la expresión aquí
        \item $\displaystyle $
    \end{enumerate}
\end{enumerate}

% --- DESAFÍO ESPECIAL (opcional, borra si no aplica) ---
\vspace{0.7cm}
\begin{desafiobox}{DESAFÍO: Título del Problema Aplicado}
\small
\textbf{Contexto:} Describe la situación real o el problema motivador.

\begin{itemize}[leftmargin=*]
    \item \textbf{Dato 1:} descripción.
    \item \textbf{Dato 2:} descripción.
\end{itemize}

\textbf{Pregunta:} ¿Qué se debe calcular o demostrar?

% Diagrama TikZ opcional:
% \begin{center}
%   \begin{tikzpicture}[scale=1.2, thick]
%     % ... tu diagrama aquí ...
%   \end{tikzpicture}
% \end{center}

\hrule\vspace{2mm}
\centering\footnotesize \textbf{Nota del Staff:} Mensaje motivacional.
\end{desafiobox}

\end{ejercicios}

% ============================================================
% SOLUCIONARIO
% ============================================================
\begin{soluciones}

\begin{solbox}{I. Soluciones — Primera Sección}
\begin{multicols}{5}
\begin{enumerate}[label=\color{azulFuerte}\bfseries\arabic*., leftmargin=*, noitemsep]
    \item $R_1$ \item $R_2$ \item $R_3$ \item $R_4$
\end{enumerate}
\end{multicols}
\end{solbox}

\begin{solbox}{II. Soluciones — Segunda Sección}
\begin{multicols}{5}
\begin{enumerate}[label=\color{azulFuerte}\bfseries\arabic*., leftmargin=*, noitemsep, start=5]
    \item $R_5$ \item $R_6$
\end{enumerate}
\end{multicols}
\end{solbox}

\begin{solbox}{III. Soluciones — Tercera Sección}
\begin{enumerate}[label=\arabic*., itemsep=1em]
    \item
    \begin{enumerate}[label=\arabic*), itemsep=0.5em]
        \item Solución a
        \item Solución b
    \end{enumerate}
\end{enumerate}
\end{solbox}

\vspace{0.8cm}
\begin{center}
    \small\textbf{Usach Premium} — ¡Nos vemos en la ayudantía! \\
    \textit{"Por y para estudiantes."}
\end{center}

\end{soluciones}

\end{document}
